\section{Cote 136: a precise uniform nonlinear test}
\label{sec:136masseytransport}

\paragraph{Scope and archive.}
The notation is the integral admissible category on physical
pp.~117 and 119 (footers 116 and 118), and the connectedness
question is physical p.~121 (footer 120),
\texttt{ARCHIVE\_136\_scan121.tex}. The result below gives a
necessary condition for integral connectedness. It does not
exhibit two theories violating it. The comparisons used are
AQ-136-COMPARISON and AQ-136-MULTIPLICATIVE; their statements
are inherited inputs. The triangular-matrix lifting below is
the usual invariance argument for degree-one defining systems;
the new point here is its single weight cutoff for all lengths.

\paragraph{Definition of the test.}
For $C\in\mathcal A_{\mathbb Z}$, a simplicial set $X$, integers
$m\geq2$ and $w_1,\ldots,w_m\geq1$, and classes
$u_i\in H^1(C^{w_i,*}(X))$, use the following convention for
the ordinary weighted Massey-product set. A defining system
is a strictly upper triangular $(m+1)$-by-$(m+1)$ matrix $A$
whose entries have form degree one and with
\[
 A_{ij}\in C^{w_i+\cdots+w_{j-1},1}(X),\qquad
 A_{1,m+1}=0,
 \qquad dA+A^2=\Omega E_{1,m+1}.                         \tag{136.MT1}
\]
The adjacent entry $A_{i,i+1}$ represents $u_i$. The element
$\Omega\in C^{w_1+\cdots+w_m,2}(X)$ is closed; its class
ranges over the Massey-product set, which may be empty.
Changing the common sign convention merely changes this
definition by its corresponding conventional signs. Nothing
here uses a selected single value of a multivalued product.

\paragraph{ARCH-136-MASSEY-TRANSPORT.}
Let $f:C\to D$ be a strict arrow of admissible integral
theories. There is $N\geq1$ such that, for every simplicial
set $X$, every length $m\geq2$, and all $w_i\geq N$, the
canonical cohomology comparisons identify the complete
weighted degree-one Massey-product sets of $C$ and $D$.
The single bound $N$ is independent of $X$, $m$, and all
the inputs. Consequently, for any two theories joined by
a finite zigzag, this equality holds with one bound for
all lengths. An inequality of these sets in arbitrarily
large input weights would therefore be a valid obstruction
to connectedness.

\begin{proof}
Choose $N$ so that AQ-136-COMPARISON identifies $H^1$ and
$H^2$ of both theories with the canonical ordinary
cohomology groups in every weight at least $N$, uniformly
in $X$. Naturality makes $f$ the identity under these
identifications. Every interval sum $w_i+\cdots+w_{j-1}$
is at least $N$, regardless of the length. In particular,
in all needed weights $H^1(f)$ is surjective and $H^2(f)$
is injective (in fact both are isomorphisms).

Applying $f$ to a defining system proves inclusion of
the $C$ set in the $D$ set. We prove the reverse inclusion,
including preservation of the entire set of possible
values. Fix a defining system $B$ in $D$, and write
$F_B=dB+B^2=\omega E_{1,m+1}$. For a strictly upper
triangular degree-zero matrix $H$ with the same interval
weights, $g=1+H$ is invertible by its finite geometric
series. The gauge transformation
\[
                  B^g=g^{-1}Bg+g^{-1}dg
\]
satisfies $F_{B^g}=g^{-1}F_Bg$. The matrix $E_{1,m+1}$
annihilates every strictly upper triangular matrix on
both sides, so $F_{B^g}=\omega E_{1,m+1}$. A gauge
transformation can create a nonzero corner entry
$B^g_{1,m+1}$. Resetting this entry to zero changes the
corner curvature by $-dB^g_{1,m+1}$ and leaves every
other entry of the curvature unchanged. It therefore
leaves the Massey value $[\omega]$ unchanged.

Construct the entries of $A$ in increasing distance
$r=j-i$, while gauging $B$ so that all entries at
already treated distances equal the corresponding
entries of $f(A)$. Only distances $1\leq r<m$ need
to be treated. At distance one choose a closed
representative of each prescribed input in $C$.
Its image and $B_{i,i+1}$ represent the same class,
so their difference is $dh$ for a degree-zero element
$h$ of that weight. Gauging by $1+hE_{i,i+1}$ changes
that entry by $dh$ and no entry at a smaller or equal
distance other than the selected one. Choose the sign
of $h$ to obtain equality, and reset the corner as above.

Suppose all smaller distances have been matched. For
$j-i=r<m$, let
\[
                       O_{ij}=\sum_{i<k<j}A_{ik}A_{kj}.
\]
The already satisfied equations $dA_{uv}=-O_{uv}$
imply $dO_{ij}=0$: the signed Leibniz rule makes each
three-factor term occur twice with opposite signs.
Its image is $-dB_{ij}$, because the smaller entries
of $B$ equal $f(A)$. Injectivity of $H^2(f)$ implies
that $O_{ij}$ is exact in $C$. Choose $a$ with
$da=-O_{ij}$. The difference $f(a)-B_{ij}$ is closed.
Surjectivity of $H^1(f)$ allows adding a closed
degree-one element to $a$ so that this difference
becomes exact. Set $A_{ij}$ equal to this corrected
element and use a gauge $1+hE_{ij}$ to make the
corresponding entries equal. All other entries at
distances at most $r$ remain unchanged, and the
corner reset preserves the represented Massey value.
All operations preserve interval weights: multiplication
of entries concatenates the corresponding intervals.

After treating all distances less than $m$, set the
corner of $A$ to zero. Its curvature is supported
only at that corner, and the curvature of its image
is the final corner curvature of $B$. The latter
represents the originally specified class $[\omega]$.
Thus this value lifts to the Massey set of $C$.
Closedness of the corner curvature follows either
from the same cancellation computation or from the
identity $dF_A=F_AA-AF_A$, whose right side is zero
for a corner-supported curvature. This also verifies
the closedness asserted in the definition.

We have proved equality using only the stated $H^1$
surjectivity and $H^2$ injectivity, with no flatness
assumption and no dependence on $m$. For a finite
zigzag take the maximum of these bounds; each arrow
then gives equality in either direction under the
canonical comparisons. This proves the last claim.
\end{proof}

\paragraph{A failed repeated-class test.}
For any odd cocycle $z\in C^{w,1}(X)$, all its strictly
commutative squares are zero. The matrix with adjacent
entries all equal to $z$ and all other entries zero
is therefore a defining system with zero curvature
for every length. Thus
\[
           0\in\langle[z],\ldots,[z]\rangle
                       \quad\text{at every finite length}.
\]
The same applies to proportional adjacent representatives.
Consequently a proposed distinction based on a nonzero
repeated-class $p$-fold product, meaning that zero is
absent from its set, cannot separate any two theories
in this strict category. This assertion does not say
that the full product set is the singleton zero, and
does not compare it with the restricted Massey operation
in ordinary singular cochains. A multivariable or
indeterminacy-sensitive construction is still required.
